% =============================================================================
% exercises/ejercicios_cap07.tex
% Ejercicios propuestos — Capítulo 7: Fisicoquímica en procesos atmosféricos
% =============================================================================

\section*{Ejercicios --- Secci\'on B: C\'alculos y aplicaciones}
\addcontentsline{toc}{section}{Ejercicios: C\'alculos y aplicaciones}
\label{sec:ejercicios_cap07b}

\begin{bloque_ejercicios}[Ejercicios B --- Fisicoquimica atmosferica]

\begin{ejercicios}

% ----- B1: Producto de solubilidad -----------------------------------------
\item El sulfato de bario (\ce{BaSO4}) tiene un $K_{ps} = 1.1\times10^{-10}$
a \SI{25}{\celsius}.

\begin{enumerate}[label=\alph*)]
    \item Escriba la ecuaci\'on de disoluci\'on y la expresi\'on del
          $K_{ps}$.
    \item Sea $C$ la solubilidad molar. Plantee $K_{ps}$ en t\'erminos
          de $C$ y calcule $C$ en mol/L.
    \item Convierta la solubilidad a \si{\milli\gram\per\liter}
          (masa molar \ce{BaSO4} = \SI{233}{\gram\per\mole}).
    \item Si se mezclan vol\'umenes iguales de \SI{0.010}{\mole\per\liter}
          de \ce{Ba^{2+}} y \SI{0.010}{\mole\per\liter} de \ce{SO4^{2-}},
          ?`precipita el \ce{BaSO4}? Calcule el producto i\'onico
          $Q = [\ce{Ba^{2+}}][\ce{SO4^{2-}}]$ y comp\'arelo con $K_{ps}$.
\end{enumerate}

\textit{[Respuesta:
(b) $K_{ps} = C^2 \Rightarrow C = \sqrt{1.1\times10^{-10}}
    = 1.05\times10^{-5}$~mol/L;
(c) $1.05\times10^{-5}\times233\times10^3
    = \SI{2.44}{\milli\gram\per\liter}$;
(d) despu\'es de mezclar $[\ce{Ba^{2+}}] = [\ce{SO4^{2-}}]
    = 0.005$~M; $Q = 2.5\times10^{-5} \gg K_{ps}$: precipita]}

% ----- B2: Ley de Henry y solubilidad de CO2 --------------------------------
\item La constante de Henry para el \ce{CO2} en agua a \SI{25}{\celsius}
es $H_{\ce{CO2}} = 3.4\times10^{-2}$~mol~L$^{-1}$~atm$^{-1}$.
La presi\'on parcial de \ce{CO2} en la atm\'osfera es
$P_{\ce{CO2}} = 4.1\times10^{-4}$~atm (410~ppm).

\begin{enumerate}[label=\alph*)]
    \item Calcule la concentraci\'on de \ce{CO2} disuelto en agua pura
          de lluvia usando $[\ce{CO2}\cdot\ce{H2O}] = H_{\ce{CO2}}P_{\ce{CO2}}$.
    \item El \'acido carb\'onico formado se disocia con
          $K_{a1} = 4.2\times10^{-7}$. Usando la aproximaci\'on
          $[\ce{H+}] \approx \sqrt{K_{a1}[\ce{H2CO3}]}$,
          estime $[\ce{H+}]$ y calcule el pH del agua de lluvia
          ``limpia''.
    \item Compare su resultado con el pH = 5.6 derivado en el
          cap\'itulo para $P_{\ce{CO2}} = 380\times10^{-6}$~atm.
          ?`Qu\'e efecto tiene el aumento de \ce{CO2} atmosf\'erico
          sobre el pH de la lluvia?
    \item Si la lluvia \'acida tiene pH~$<$~5.0, ?`qu\'e
          concentraci\'on de \ce{SO2} adicional (en ppm) ser\'ia
          necesaria para bajar el pH a ese nivel?
          Use $H_{\ce{SO2}} = 1.2$~M/atm y $k'_1 = 1.71\times10^{-2}$.
\end{enumerate}

\textit{[Respuesta:
(a) $[\ce{CO2}\cdot\ce{H2O}] = 3.4\times10^{-2}
    \times 4.1\times10^{-4} = 1.39\times10^{-5}$~mol/L;
(b) $[\ce{H+}] = \sqrt{4.2\times10^{-7}\times1.39\times10^{-5}}
    = 2.42\times10^{-6}$~M; pH~$= 5.62$]}

% ----- B3: Gradiente adiabático seco y estabilidad atmosférica -------------
\item El gradiente adiab\'atico seco se define como
$\Gamma = -g/C_p = \SI{9.76}{\celsius\per\kilo\metre}$.

\begin{enumerate}[label=\alph*)]
    \item Una parcela de aire a nivel del suelo tiene
          $T_0 = \SI{25}{\celsius}$. Si asciende adiab\'aticamente
          \SI{3}{\kilo\metre}, calcule su temperatura final.
    \item Si el gradiente ambiental observado es de
          \SI{7}{\celsius\per\kilo\metre} (condici\'on est\'able),
          ?`cu\'al ser\'a la temperatura del ambiente a \SI{3}{\kilo\metre}
          si $T_0^{\text{amb}} = \SI{25}{\celsius}$?
    \item Compare la temperatura de la parcela con la del ambiente
          a \SI{3}{\kilo\metre}. ?`La parcela asciende o desciende
          espont\'aneamente? Explique en t\'erminos de densidad.
    \item ?`Qu\'e ocurrir\'ia si el gradiente ambiental fuera de
          \SI{12}{\celsius\per\kilo\metre} (superadiab\'atico)?
          Describa la estabilidad de la atm\'osfera en ese caso.
\end{enumerate}

\textit{[Respuesta:
(a) $T_{\text{parcela}} = 25 - 9.76\times3 = \SI{-4.3}{\celsius}$;
(b) $T_{\text{amb}} = 25 - 7\times3 = \SI{4}{\celsius}$;
(c) parcela a \SI{-4.3}{\celsius} $<$ ambiente a \SI{4}{\celsius}:
    la parcela es m\'as densa y desciende $\Rightarrow$ atm\'osfera
    estable]}

% ----- B4: Aire húmedo y flotabilidad --------------------------------------
\item Con base en las ecuaciones de aire h\'umedo del cap\'itulo
($M_a = \SI{29}{\gram\per\mole}$, $M_{\ce{H2O}} = \SI{18}{\gram\per\mole}$):

\begin{enumerate}[label=\alph*)]
    \item Calcule la masa molecular promedio del aire h\'umedo
          $M_w = (1-f_H)M_a + f_H M_{\ce{H2O}}$ para fracciones
          h\'umedas $f_H = 0$, 0.02 y 0.04.
    \item Calcule la raz\'on $m_d/m_w = M_a/M_w$ para cada
          caso del inciso (a). ?`C\'omo cambia con $f_H$?
    \item Usando la ecuaci\'on de aceleraci\'on por flotaci\'on:
          \[
              B = \frac{(M_a - M_{\ce{H2O}})f_H}
                       {M_a - (M_a-M_{\ce{H2O}})f_H}\,g
          \]
          calcule $B$ para $f_H = 0.02$ y $f_H = 1.0$.
          (Use $g = \SI{9.81}{\metre\per\square\second}$)
    \item Explique f\'isicamente por qu\'e el aire h\'umedo
          asciende m\'as f\'acilmente que el aire seco.
\end{enumerate}

\textit{[Respuesta:
(a) $f_H=0$: $M_w=29$; $f_H=0.02$: $M_w=28.78$;
    $f_H=0.04$: $M_w=28.56$~g/mol;
(c) $f_H=0.02$: $B = (11\times0.02)/(29-11\times0.02)
    \times9.81 = 0.22/28.78\times9.81
    \approx \SI{0.075}{\metre\per\square\second}$;
    $f_H=1.0$: $B \approx \SI{6}{\metre\per\square\second}$]}

% ----- B5: Cinética de primer orden — tiempo de vida media -----------------
\item La descomposici\'on del \ce{NO2} sigue cin\'etica de primer orden
con constante $k = 3.0\times10^{-3}$~s$^{-1}$.

\begin{enumerate}[label=\alph*)]
    \item Calcule el tiempo de vida media $t_{1/2} = \ln2/k$.

    \item Si la concentraci\'on inicial de \ce{NO2} es
          $[\ce{NO2}]_0 = 200$~ppb, calcule la concentraci\'on
          restante despu\'es de 1, 3 y 5 tiempos de vida media.

    \item Grafique cualitativamente la disminuci\'on de [\ce{NO2}]
          con el tiempo, indicando los puntos calculados en (b).

    \item La fot\'olisis del \ce{NO2} con luz solar sigue la reacci\'on
          \ce{NO2 ->[$h\nu$] NO + O^.} con una constante de velocidad
          fotol\'itica $j = 8\times10^{-3}$~s$^{-1}$ al mediod\'ia.
          Calcule $t_{1/2}$ para este proceso y compare con el (a).
\end{enumerate}

\textit{[Respuesta:
(a) $t_{1/2} = 0.693/3.0\times10^{-3} = \SI{231}{\second}
    \approx \SI{3.85}{\minute}$;
(b) $n=1$: 100~ppb; $n=3$: 25~ppb; $n=5$: 6.25~ppb;
(d) $t_{1/2} = 0.693/8\times10^{-3}
    = \SI{86.6}{\second} \approx \SI{1.44}{\minute}$;
    la fotolisis es $\approx$2.7 veces m\'as r\'apida]}

% ----- B6: Tiempo de residencia atmosférico --------------------------------
\item Con base en el \textbf{Cuadro~\ref{tresidencia}}
de tiempos de residencia:

\begin{enumerate}[label=\alph*)]
    \item Complete la siguiente tabla indicando la escala espacial
          dominante (local, regional o global) que corresponde
          a cada gas seg\'un su tiempo de residencia:

{%
\centering
\begin{tabular}{lcl}
\toprule
\textbf{Gas} & \textbf{Tiempo de residencia} & \textbf{Escala} \\
\midrule
\ce{NO2}  & 0.5--2 d\'ias  & \\
CO        & $\sim$60 d\'ias & \\
\ce{O3}   & $\sim$100 d\'ias & \\
\ce{CH4}  & 9 a\~nos        & \\
\ce{CO2}  & 3--4 a\~nos     & \\
\ce{N2}   & $1.6\times10^7$ a\~nos & \\
\bottomrule
\end{tabular}
\par
}%

    \item La masa total de \ce{CH4} en la atm\'osfera es
          $M \approx 5\times10^{12}$~kg. Usando $\tau = M/F$,
          calcule la tasa de emisi\'on anual $F$ en kg/a\~no
          (\textit{use} $\tau = 9$~a\~nos).
    \item Explique por qu\'e el \ce{NO2} se acumula cerca de
          fuentes urbanas pero no contribuye a la contaminaci\'on
          de fondo global, mientras que el \ce{CH4} s\'i.
\end{enumerate}

\textit{[Respuesta:
(b) $F = M/\tau = 5\times10^{12}/9
    \approx 5.6\times10^{11}$~kg/a\~no
    $= \SI{560e9}{\kilo\gram}/$a\~no]}

% ----- B7: Cinética de segundo orden — reacción NO + O3 -------------------
\item La reacci\'on \ce{O3 + NO -> NO2 + O2} es de segundo orden con
$k = 1.9\times10^{-14}$~cm$^3$~mol\'ec$^{-1}$~s$^{-1}$ a
\SI{298}{\kelvin}.

\begin{enumerate}[label=\alph*)]
    \item Escriba la expresi\'on de la ley de velocidad y sus unidades.
    \item En una masa de aire urbana con
          $[\ce{O3}] = 100$~ppb $= 2.46\times10^{12}$~mol\'ec/cm$^3$
          y $[\text{NO}] = 50$~ppb $= 1.23\times10^{12}$~mol\'ec/cm$^3$,
          calcule la velocidad de la reacci\'on en
          mol\'ec~cm$^{-3}$~s$^{-1}$.
    \item Si $[\ce{O3}]$ est\'a en exceso (pseudo-primer orden),
          calcule la constante efectiva $k' = k[\ce{O3}]$ y el
          $t_{1/2}$ del NO.
    \item Esta reacci\'on es clave en el ciclo de smog fotoqu\'imico.
          Explique por qu\'e su velocidad aumenta en el mediod\'ia
          cuando [\ce{O3}] es m\'axima.
\end{enumerate}

\textit{[Respuesta:
(b) $v = 1.9\times10^{-14}
    \times 2.46\times10^{12}\times1.23\times10^{12}
    \approx 5.75\times10^{10}$~mol\'ec~cm$^{-3}$~s$^{-1}$;
(c) $k' = 1.9\times10^{-14}\times2.46\times10^{12}
    = 4.67\times10^{-2}$~s$^{-1}$;
    $t_{1/2} = 0.693/4.67\times10^{-2}
    \approx \SI{14.8}{\second}$]}

% ----- B8: Fotoquímica y fotodisociación -----------------------------------
\item La fotodisociaci\'on del \ce{NO2} es el paso iniciador del
smog fotoqu\'imico: \ce{NO2 ->[$j$] NO + O^.}, donde
$j$ es la constante de fotolisis.

\begin{enumerate}[label=\alph*)]
    \item La energ\'ia de un fot\'on es $E = h\nu = hc/\lambda$.
          Para $\lambda = \SI{400}{\nano\metre}$, calcule $E$ en
          joules y en eV.
          (Datos: $h = 6.626\times10^{-34}$~J~s,
          $c = 3\times10^8$~m/s, 1~eV = $1.602\times10^{-19}$~J)
    \item La energ\'ia de enlace \ce{N-O} en \ce{NO2} es
          $\approx$~\SI{305}{\kilo\joule\per\mole}.
          Exprese este valor en eV/mol\'ecula y determine si la
          luz visible ($\lambda \leq$ \SI{420}{\nano\metre}) tiene
          energ\'ia suficiente para disociar el \ce{NO2}.
    \item Describa los dos pasos de la fotodisociaci\'on de
          una mol\'ecula A: absorci\'on del fot\'on y disociaci\'on
          del estado excitado, con las ecuaciones indicadas
          en el cap\'itulo.
    \item Explique por qu\'e la formaci\'on de smog fotoqu\'imico
          es m\'axima al mediod\'ia y por qu\'e no ocurre de noche.
\end{enumerate}

\textit{[Respuesta:
(a) $E = 6.626\times10^{-34}\times3\times10^8
    / 4\times10^{-7}
    = 4.97\times10^{-19}$~J $= 3.10$~eV;
(b) $305000/(6.022\times10^{23}\times1.602\times10^{-19})
    = 3.16$~eV/mol\'ecula;
    la luz a 400~nm (3.10~eV) est\'a justo en el l\'imite:
    la luz UV cercana ($\lambda < 400$~nm) s\'i dissocia]}

% ----- B9: Visibilidad y ley de Beer-Lambert --------------------------------
\item La ley de Beer-Lambert extendida describe la atenuaci\'on
de intensidad luminosa:
\[
    I = I_0\,e^{-(\alpha+s)d}
\]
donde $\alpha$ es el coeficiente de absorci\'on,
$s$ el de esparcimiento y $d$ la distancia.

\begin{enumerate}[label=\alph*)]
    \item Si $\alpha + s = \beta_{ext} = 200$~Mm$^{-1}$
          $= 2\times10^{-4}$~m$^{-1}$, calcule la fracci\'on
          de intensidad $I/I_0$ que llega al observador
          a $d = 10$~km y a $d = 20$~km.
    \item Usando la f\'ormula de Koschmieder
          $V = 3.912/\beta_{ext}$, calcule la visibilidad
          para $\beta_{ext} = 200$~Mm$^{-1}$ y
          $\beta_{ext} = 39\,000$~Mm$^{-1}$ (niebla densa).
    \item La dispersi\'on de Mie es m\'axima para part\'iculas
          con radio $r \sim \lambda/2\pi$. Para la luz verde
          ($\lambda = \SI{0.52}{\micro\metre}$), calcule el
          radio de part\'icula que maximiza la dispersi\'on.
    \item Seg\'un el cap\'itulo, las part\'iculas de
          \SIrange{0.1}{1}{\micro\metre} son las m\'as efectivas
          en reducir la visibilidad. ?`En qu\'e rango de
          PM caen estas part\'iculas (PM$_{10}$ o PM$_{2.5}$)?
          ?`Por qu\'e tambi\'en afectan a la salud?
\end{enumerate}

\textit{[Respuesta:
(a) $I/I_0 = e^{-2\times10^{-4}\times10000}
    = e^{-2} = 0.135$ (13.5\%) a 10~km;
    $e^{-4} = 0.018$ (1.8\%) a 20~km;
(b) $V_{200} = 3.912/0.200 = \SI{19.6}{\kilo\metre}$;
    $V_{39000} = 3.912/39.0 \approx \SI{0.10}{\kilo\metre}$;
(c) $r = \lambda/(2\pi)
    = 0.52/(2\pi) \approx \SI{0.083}{\micro\metre}$]}

% ----- B10 (avanzado): Equilibrio global de la lluvia ácida ---------------
\item \textbf{[Avanzado]} En una masa de aire contaminada se tienen
las siguientes presiones parciales:
$P_{\ce{CO2}} = 380\times10^{-6}$~atm,
$P_{\ce{SO2}} = 10\times10^{-9}$~atm (10~ppb) y
$P_{\ce{NH3}} = 2\times10^{-9}$~atm (2~ppb).

Usando los par\'ametros del cap\'itulo:
$H_{\ce{CO2}} = K_H(\ce{CO2})$, $k_1 = 4.2\times10^{-7}$,
$H_{\ce{SO2}} = 1.2$~M/atm, $k'_1 = 1.71\times10^{-2}$,
$H_{\ce{NH3}} = 58$~M/atm, $k_b = 5.6\times10^{-10}$,
$K_w = 10^{-14}$:

\begin{enumerate}[label=\alph*)]
    \item Calcule las concentraciones acuosas iniciales:
          $[\ce{CO2}\cdot\ce{H2O}]$, $[\ce{SO2}\cdot\ce{H2O}]$
          y [\ce{NH3}$_{(ac)}$] usando la ley de Henry.
    \item ?`Cu\'al de los tres gases tiene mayor concentraci\'on
          disuelta? ?`Cu\'al es el que m\'as acidifica y cu\'al
          el que m\'as neutraliza la lluvia?
    \item Para estimar el efecto del \ce{SO2} solo, calcule
          $[\ce{HSO3^-}]$ suponiendo $[\ce{H^+}] = 10^{-5.6}$~M
          (el pH del agua limpia de lluvia) como primera
          aproximaci\'on.
    \item Comente cualitativamente si el pH resultante ser\'a
          mayor o menor que 5.6 considerando los tres gases
          simult\'aneamente, y justifique con base en las
          constantes de equilibrio.
\end{enumerate}

\textit{[Respuesta parcial:
(a) $[\ce{CO2}\cdot\ce{H2O}]
    = H_{\ce{CO2}}\times380\times10^{-6}
    \approx 1.3\times10^{-5}$~M;
    $[\ce{SO2}\cdot\ce{H2O}]
    = 1.2\times10\times10^{-9} = 1.2\times10^{-8}$~M;
   [\ce{NH3} $_{(ac)}$] $= 58\times2\times10^{-9}
    = 1.16\times10^{-7}$~M;
(b) \ce{CO2} tiene mayor concentraci\'on;
    \ce{SO2} acidifica m\'as (k$'_1 = 1.71\times10^{-2}
    \gg k_1 = 4.2\times10^{-7}$); \ce{NH3} neutraliza]}

\end{ejercicios}

\end{bloque_ejercicios}
% =============================================================================
% FIN DE ejercicios_cap07.tex
% =============================================================================
